% =====================================================================
%  Lecture 16 -- Network Functions Virtualisation (NFV)   (English edition)
% =====================================================================
\documentclass[11pt,a4paper]{article}
\providecommand\ACNroot{../..}
\input{\ACNroot/templates/lecture-en}

\ACNset{number}{16}
\ACNset{title}{Network Functions Virtualisation (NFV)}
\ACNset{shorttitle}{NFV}
\ACNset{titleother}{مجازی‌سازی کارکردهای شبکه (NFV)}
\ACNset{part}{Part III -- Software-driven architectures (Lectures 15--18)}
\ACNset{time}{6--8 hours (3 h reading, 3--5 h exercises)}
\ACNset{level}{Intermediate}
\ACNset{prereq}{Lectures 11 and 15; operating systems (processes, virtual memory); basic probability}
\ACNset{gen}{5}

\begin{document}
\ACNcover

% ---------------------------------------------------------------------
\section{Learning objectives}
After studying this lecture you should be able to:
\begin{objectives}
  \item \textbf{Explain} why operators moved network functions from dedicated appliances to software on
        shared infrastructure.
  \item \textbf{Describe} the ETSI NFV architectural framework: VNFs, NFVI and the MANO functions with their
        reference points.
  \item \textbf{Compare} virtual machines and containers, and \textbf{explain} cloud-native network functions.
  \item \textbf{Trace} the VNF lifecycle, including scaling and healing.
  \item \textbf{Compute} availability of virtualised stacks, scaling needs and per-packet processing budgets.
  \item \textbf{Relate} NFV to SDN and to the virtualised EPC and the cloud-native 5G core.
\end{objectives}

% ---------------------------------------------------------------------
\section{Prerequisite knowledge}
You should know what an operating system kernel, a process and virtual memory are, and how series and
parallel availability are combined (Lecture~11). SDN (Lecture~15) is the natural companion of NFV.

\paragraph{Key terms} The vocabulary of this lecture with its Persian equivalents:
\begin{keyterms}
NFV & \fa{مجازی‌سازی کارکردهای شبکه} & Running network functions as software on general-purpose servers. \\
VNF / VNFC & \fa{کارکرد شبکهٔ مجازی / مؤلفه} & A virtualised network function / one of its components. \\
NFVI & \fa{زیرساخت NFV} & Compute, storage, network hardware and the virtualisation layer. \\
MANO & \fa{مدیریت و هماهنگ‌سازی} & NFV orchestrator, VNF manager and virtualised infrastructure manager. \\
VNFD / NSD & \fa{توصیفگر VNF / خدمت شبکه} & Templates describing a VNF / a network service. \\
Hypervisor & \fa{ابرناظر} & Software that runs virtual machines on a host. \\
Container & \fa{کانتینر} & Isolated process group sharing the host kernel. \\
CNF & \fa{کارکرد شبکهٔ ابرزاد} & Cloud-native network function built from containers/microservices. \\
SR-IOV / DPDK & \fa{فناوری‌های شتاب I/O} & Direct NIC access for VMs / user-space packet processing. \\
Service function chain & \fa{زنجیرهٔ کارکرد خدمت} & Ordered sequence of functions a flow traverses. \\
\end{keyterms}

% ---------------------------------------------------------------------
\section{Introduction and motivation}
Until the 2010s every network function -- MME, HSS, gateway, firewall, load balancer -- was a dedicated
box: specialised hardware, a proprietary operating system and the function software, sold together.
Adding capacity meant buying and installing more boxes; introducing a service meant integrating yet another
appliance; each type had its own spares, tools and lifecycle. In 2012 a group of operators proposed, in a
white paper that launched the ETSI NFV Industry Specification Group, to run these functions as software
on standard servers in data centres, managed by common orchestration \cite{etsinfv002}.

The promised benefits are those of cloud computing: shared hardware and lower cost, capacity that follows
demand (elastic scaling), faster introduction of services, and automation. The costs are equally real:
packet processing on general-purpose CPUs is harder than on ASICs, and availability now depends on a stack
of hardware, virtualisation and software layers operated by different teams. The virtualised EPC was the
first large mobile application; the 5G core was designed from the start as cloud-native software
(Lectures~18--19).

% ---------------------------------------------------------------------
\section{Fundamental concepts}

\subsection{The ETSI NFV framework}
The framework (\cref{fig:framework}) \cite{etsinfv002} has three domains:
\begin{itemize}
  \item \textbf{VNFs}, each with an element manager (EM) for its application-level configuration.
  \item \textbf{NFVI}: physical compute, storage and network resources, abstracted by a virtualisation layer
        into virtual resources.
  \item \textbf{MANO} \cite{etsinfvman001}: the \emph{NFV orchestrator} (NFVO) manages network services and
        global resources; the \emph{VNF manager} (VNFM) manages the lifecycle of VNF instances; the
        \emph{virtualised infrastructure manager} (VIM, e.g.\ an OpenStack or Kubernetes control plane)
        manages the NFVI resources.
\end{itemize}
Named reference points (Os-Ma, Or-Vnfm, Vi-Vnfm, Or-Vi, Ve-Vnfm, Nf-Vi, Vn-Nf, Vl-Ha) define who talks to whom,
so that components from different vendors can be combined.

\begin{acnfigure}
  \resizebox{\linewidth}{!}{\acnfig{lecture-16/nfv-framework}}
  \caption{ETSI NFV architectural framework with its main reference points \cite{etsinfv002}.}
  \label{fig:framework}
\end{acnfigure}

\subsection{Descriptors and network services}
A VNF is delivered as a package containing software images and a \emph{VNF descriptor} (VNFD): its components
(VNFCs), their resource needs, internal virtual links, connection points, scaling aspects and lifecycle
scripts. A \emph{network service} combines VNFs and the virtual links between them, described by a network
service descriptor (NSD) and, where traffic must traverse functions in order, a forwarding graph. Descriptors
make deployment declarative: the orchestrator turns a template into running instances.

\subsection{Virtual machines and containers}
A \emph{virtual machine} runs a complete guest operating system on a hypervisor; it gives strong isolation and
supports any OS, at the cost of memory and boot time. A \emph{container} isolates a group of processes with kernel
namespaces and control groups while sharing the host kernel; it starts in seconds and has little overhead
(\cref{fig:stacks}). \emph{Cloud-native network functions} (CNFs) are built as sets of containerised
microservices, keep session state in separate data stores so that instances are disposable, and are deployed,
scaled and healed by a container orchestrator such as Kubernetes.

\begin{acnfigure}
  \resizebox{\linewidth}{!}{\acnfig{lecture-16/virt-stacks}}
  \caption{From appliances to virtual machines and containers.}
  \label{fig:stacks}
\end{acnfigure}

% ---------------------------------------------------------------------
\section{Detailed technical explanation}

\subsection{Lifecycle management}
The VNFM drives a VNF through the lifecycle operations of the MANO specifications \cite{etsinfvman001}:
\emph{onboarding} (the package is validated and its images stored), \emph{instantiation}, \emph{scaling} out/in
(adding or removing VNFC instances) or up/down (changing their resources), \emph{healing} (replacing failed
instances), \emph{update/upgrade}, and \emph{termination}. Scaling and healing can be triggered by the VNF or EM,
by the VNFM from performance and fault data, or by the NFVO for a whole network service (\cref{fig:scale}).

\subsection{High-performance data planes}
Gateways and UPFs must forward millions of packets per second. Plain virtual switching through the kernel and
hypervisor cannot keep up, so NFVI nodes for user-plane functions use:
\begin{itemize}
  \item \textbf{DPDK}: polling-mode drivers and packet processing in user space, avoiding interrupts and kernel
        copies;
  \item \textbf{SR-IOV}: the NIC presents virtual functions that are passed directly to VMs or containers;
  \item \textbf{CPU pinning and NUMA awareness}: dedicated cores on the socket attached to the NIC, huge pages for
        packet buffers;
  \item \textbf{Hardware offload}: smart NICs or programmable switches (Lecture~15) for encapsulation, classification
        or encryption.
\end{itemize}
These techniques trade some of NFV's flexibility (portability, overcommitment) for performance; the per-packet
budget of \cref{sec:cycles} explains why they are needed.

\subsection{Resilience in a virtualised world}
In an appliance, redundancy was built into the box (redundant boards, synchronised state). In NFV, resilience is a
property of the software design: instances are spread across hosts and sites (anti-affinity rules), state is
replicated or externalised, and failed instances are healed by the orchestrator. Common-mode failures move up the
stack: a faulty hypervisor update or orchestrator bug can affect every function at once.

\subsection{NFV in mobile networks}
The virtualised EPC (vEPC) packaged MME, S/P-GW, HSS and PCRF as VNFs, initially as ``lift and shift'' of
appliance software into VMs. The 5G core takes the next step: network functions are defined as services
(Lecture~18), implemented as CNFs and deployed on Kubernetes-based platforms, with the user plane at edge sites.
In the RAN, virtualised and Cloud RAN run the baseband (CU and DU) on servers with acceleration hardware.

% ---------------------------------------------------------------------
\section{Mathematical models}

\subsection{Availability of a virtualised stack}\label{sec:avail}
\paragraph{Model} A VNF instance needs its host hardware, the virtualisation layer and its own software, with
availabilities $a_{\mathrm{hw}}$, $a_{\mathrm{v}}$ and $a_{\mathrm{sw}}$ (independent). In series, and with $k$
independent redundant instances of which one suffices,
\begin{equation}\label{eq:series}
  A_1 = a_{\mathrm{hw}}\,a_{\mathrm{v}}\,a_{\mathrm{sw}}, \qquad
  A_{k} = 1-(1-A_1)^{k}.
\end{equation}
\paragraph{Example} $a_{\mathrm{hw}} = 0.999$, $a_{\mathrm{v}} = 0.9995$, $a_{\mathrm{sw}} = 0.999$: $A_1 = 0.9975$, i.e.\ 21.9\,h
of downtime per year -- worse than each component. Two instances on independent hosts: $A_2 = 1 - (0.0025)^2 =
0.999994$, about 197\,s per year, provided failover is instantaneous.
\paragraph{Interpretation} Redundancy at the instance level more than compensates for the less reliable layers,
but only if failures are independent (different hosts, racks, sites) and state survives the failover.

\subsection{Elastic scaling}\label{sec:scale}
\paragraph{Model} An instance serves $\mu$ requests/s; the operator keeps utilisation at most $\rho_{\mathrm{t}}$. For
offered load $\lambda$ the number of instances is
\begin{equation}\label{eq:scale}
  n = \left\lceil \frac{\lambda}{\rho_{\mathrm{t}}\,\mu} \right\rceil .
\end{equation}
If load can grow at rate $g$ (requests/s per second) and a new instance takes $T_\mathrm{i}$ to become ready, the
spare capacity must cover the growth during instantiation:
\begin{equation}\label{eq:headroom}
  n\mu - \lambda \ge g\,T_\mathrm{i}
  \quad\Longleftrightarrow\quad
  \rho \le 1 - \frac{g\,T_\mathrm{i}}{n\mu}.
\end{equation}
\paragraph{Example} $\lambda = 40\,000$\,s$^{-1}$, $\mu = 5000$\,s$^{-1}$, $\rho_{\mathrm{t}} = 0.7$: $n = \lceil 11.4\rceil = 12$. With
$g = 100$\,s$^{-2}$ and $T_\mathrm{i} = 60$\,s the headroom must be 6000\,s$^{-1}$, so $\rho \le 1 - 6000/60\,000 = 0.9$: the
70\,\% target is safe. Containers ($T_\mathrm{i}$ of seconds) allow much higher utilisation than VMs (minutes).

\subsection{Per-packet processing budget}\label{sec:cycles}
\paragraph{Model} At line rate $R$ with packets of $L$ bytes (plus 20 bytes of Ethernet preamble and inter-frame
gap), one core at clock $f$ has
\begin{equation}\label{eq:cycles}
  c = \frac{f}{R/\bigl(8(L+20)\bigr)}
\end{equation}
cycles per packet.
\paragraph{Example} $f = 2.5$\,GHz, $R = 10$\,Gbit/s, $L = 64$: $14.88$\,Mpps and $c = 168$ cycles -- less than the cost
of a single cache miss to main memory in the worst case. At 100\,Gbit/s only 17 cycles remain (\cref{fig:cycles}).
\paragraph{Interpretation} Line-rate small-packet forwarding needs many cores in parallel, batching, cache-friendly
data structures and offload; real traffic with larger average packets relaxes the budget considerably.

\begin{acnfigure}
  \acnfig{lecture-16/cycles}
  \caption{Cycles available per packet on one 2.5\,GHz core at line rate (\cref{eq:cycles}).}
  \label{fig:cycles}
\end{acnfigure}

\subsection{Placement as bin packing}\label{sec:pack}
\paragraph{Model} VNFCs with CPU demands $d_j$ must be placed on servers of capacity $C$. The number of servers is at
least
\begin{equation}\label{eq:binpack}
  m \ge \left\lceil \frac{\sum_j d_j}{C} \right\rceil ,
\end{equation}
and finding the minimum is the NP-hard bin-packing problem; heuristics such as first-fit decreasing come close.
Anti-affinity (replicas on different servers) and latency constraints (user plane at the edge) add constraints that
can increase $m$.

% ---------------------------------------------------------------------
\section{Network architecture and diagrams}
An operator's NFV deployment typically consists of a few central data centres (control-plane functions,
subscriber databases, OSS) and many edge sites (user-plane functions, RAN baseband), all managed by a common
orchestration layer. SDN controllers (Lecture~15) configure the data-centre fabrics and the transport between sites,
so that the virtual links of network services are realised end to end. \Cref{fig:framework} shows the functional
view; physically, MANO components themselves run as virtualised applications on the same infrastructure.

% ---------------------------------------------------------------------
\section{Step-by-step procedures}

\begin{procedurebox}{Instantiating a network service}\label{proc:inst}
\begin{steps}
  \item The operator (OSS) requests instantiation of a network service from the NFVO, referencing an onboarded NSD.
  \item The NFVO checks resources, selects VIMs/sites and asks the VNFM of each VNF to instantiate it.
  \item The VNFM requests a resource grant from the NFVO, then asks the VIM to create VMs or containers, networks and
        storage according to the VNFD.
  \item The VIM allocates the resources on the NFVI and boots the images; the VNFM configures the instances and
        informs the EM.
  \item The NFVO creates the virtual links between the VNFs (possibly through an SDN controller) and reports the
        service as instantiated.
\end{steps}
\end{procedurebox}

\begin{procedurebox}{Automatic scale-out}\label{proc:scale}
\begin{steps}
  \item The VNF or EM reports load above a threshold defined in the VNFD (or the VNFM detects it from metrics).
  \item The VNFM decides the scaling step and obtains a grant from the NFVO.
  \item The VIM allocates resources and starts a new VNFC instance.
  \item The VNFM configures it and adds it to the VNF's internal load distribution; traffic shifts to it.
\end{steps}
\end{procedurebox}

\begin{acnfigure}
  \resizebox{\linewidth}{!}{\acnfig{lecture-16/scale-out}}
  \caption{Simplified VNF scale-out (\cref{proc:scale}).}
  \label{fig:scale}
\end{acnfigure}

% ---------------------------------------------------------------------
\section{Worked examples}

\begin{examplebox}{Redundancy design}\label{ex:avail}
\textbf{Problem.} With the component availabilities of \cref{sec:avail}, compare one instance, two instances on the same
host, and two instances on different hosts.

\textbf{Solution.} One: $A_1 = 0.9975$ (21.9\,h/year). Same host: hardware and hypervisor are shared, so
$A = a_{\mathrm{hw}}a_{\mathrm{v}}\,[1-(1-a_{\mathrm{sw}})^2] = 0.9985$ (13.1\,h/year): only software failures are covered.
Different hosts: $A_2 = 0.999994$ (197\,s/year). Anti-affinity is essential.
\end{examplebox}

\begin{examplebox}{Scaling a signalling function}\label{ex:scale}
\textbf{Problem.} An AMF-like function must handle a busy hour of $40\,000$ procedures/s; each instance handles 5000/s.
How many instances are needed at 70\,\% target utilisation, and what is the utilisation after one instance fails?

\textbf{Solution.} $n = 12$ (\cref{eq:scale}); utilisation $40\,000/60\,000 = 67\,\%$. After a failure: $40\,000/55\,000 = 73\,\%$ --
still below saturation, and the orchestrator heals the instance in the meantime.
\end{examplebox}

\begin{examplebox}{Packing VNFCs}\label{ex:pack}
\textbf{Problem.} VNFCs need 6, 5, 5, 4, 4, 3, 3 and 2 cores; servers have 16 cores. Find the lower bound and a first-fit
decreasing placement.

\textbf{Solution.} $\sum d_j = 32$, so $m \ge 2$. First-fit decreasing: server 1 gets 6, 5, 5 (16 cores); server 2 gets
4, 4, 3, 3, 2 (16 cores). Two servers, which is optimal -- but with no spare capacity for healing.
\end{examplebox}

% ---------------------------------------------------------------------
\section{Technology comparison}
\begin{table}[htbp]
\centering\small
\caption{Physical appliances, VM-based VNFs and cloud-native CNFs.}\label{tab:cmp}
\rowcolors{2}{acnmist}{white}
\begin{tabularx}{\linewidth}{@{}l Y Y Y@{}}
\toprule
\rowcolor{white} & Appliance & VNF (VMs) & CNF (containers) \\
\midrule
Hardware          & Dedicated, vendor-specific & Shared COTS servers & Shared COTS servers \\
Isolation         & Physical                   & Hypervisor          & Kernel namespaces, cgroups \\
Start-up time     & Hours to days (install)    & Minutes             & Seconds \\
Scaling           & Buy and install boxes      & Add VMs (scale out/up) & Add pods automatically \\
State handling    & Inside the box             & Often inside the VM & Externalised data stores \\
Lifecycle         & Vendor-specific tools      & MANO (NFVO, VNFM, VIM) & Kubernetes + operators, CI/CD \\
Data-plane speed  & ASIC line rate             & Needs DPDK, SR-IOV  & Needs DPDK, SR-IOV, offload \\
\bottomrule
\end{tabularx}
\end{table}

\paragraph{NFV and SDN} The two are complementary and independent: NFV changes \emph{where} functions run (software on
shared servers), SDN changes \emph{how} forwarding is controlled (programmable devices, central control). Together
they allow a network service -- functions and the links between them -- to be created by software.

% ---------------------------------------------------------------------
\section{Key takeaways}
\begin{takeawaybox}
\begin{itemize}
  \item NFV runs network functions as software on shared infrastructure, managed by MANO (NFVO, VNFM, VIM).
  \item Descriptors (VNFD, NSD) make deployment declarative; lifecycle operations include scale and heal.
  \item Containers and cloud-native design (stateless instances, external state) underpin the 5G core.
  \item Series availability drops with each layer; independent redundant instances restore it: $A_k = 1-(1-A_1)^k$.
  \item Scaling: $n = \lceil\lambda/(\rho_{\mathrm{t}}\mu)\rceil$ with headroom $gT_\mathrm{i}$; packet budget $c = 8f(L+20)/R$ cycles.
\end{itemize}
\end{takeawaybox}

% ---------------------------------------------------------------------
\section{Exercises}

\begin{exercisebox}{1}{Who does what?}\label{exr:mano}
Which MANO component (NFVO, VNFM, VIM) is responsible for: (a) creating a virtual machine; (b) deciding to add a
VNFC instance; (c) connecting several VNFs into a network service; (d) keeping an inventory of NFVI resources?
\end{exercisebox}

\begin{exercisebox}{1}{VM or container?}\label{exr:vmct}
Give one argument for running a virtualised HSS in VMs and one for running it as containers.
\end{exercisebox}

\begin{exercisebox}{2}{Same-host redundancy}\label{exr:host}
Using \cref{eq:series}, compute the annual downtime of three instances placed on three different hosts that share one
power and network domain with availability 0.9999. Which term dominates?
\end{exercisebox}

\begin{exercisebox}{2}{Packet budget}\label{exr:cycles}
A UPF core runs at 3\,GHz and needs 400 cycles per packet. What throughput can one core sustain with 64-byte and with
800-byte packets? How many cores for 100\,Gbit/s at 800 bytes?
\end{exercisebox}

\begin{exercisebox}{3}{Scaling lag}\label{exr:lag}
Load ramps at $g = 200$\,s$^{-2}$. Instances serve $\mu = 4000$\,s$^{-1}$. Compare the maximum safe utilisation for
$n = 10$ instances when a new instance takes 120\,s (VM) or 10\,s (container) to become ready. Discuss what else limits
the speed of scaling.
\end{exercisebox}

\begin{exercisebox}{3}{Placement with anti-affinity}\label{exr:pack}
The VNFCs of \cref{ex:pack} consist of two replicas of a 6-core component (A), two of a 5-core component (B) and the rest
(4, 4, 3, 3, 2) without constraints; replicas of the same component must be on different servers, and each server must
keep 2 cores free for healing. Find a placement with the fewest 16-core servers.
\end{exercisebox}

\subsection*{Solutions}
\addcontentsline{toc}{subsection}{Solutions}

\begin{solutionbox}{\ref{exr:mano}}
(a) VIM; (b) VNFM (possibly at the request of the EM/VNF or NFVO); (c) NFVO; (d) VIM (with a global view kept by the NFVO).
\end{solutionbox}

\begin{solutionbox}{\ref{exr:vmct}}
VMs: strong isolation for a security-critical database holding subscriber keys, and support for an existing vendor image.
Containers: fast scaling and healing, smaller footprint and integration with cloud-native tooling, if the HSS software
is designed with external, replicated state.
\end{solutionbox}

\begin{solutionbox}{\ref{exr:host}}
Three independent instances: $1-(0.0025)^3 \approx 1-1.6\times10^{-8}$. Shared domain in series: $A = 0.9999\times(1-1.6\times10^{-8})
\approx 0.99990$, i.e.\ about 53\,min/year, almost all due to the shared domain. Independence of failure domains matters more
than the number of instances.
\end{solutionbox}

\begin{solutionbox}{\ref{exr:cycles}}
One core: $3\times10^9/400 = 7.5$\,Mpps. 64-byte packets: $7.5\times10^6\times84\times8 = 5.0$\,Gbit/s; 800-byte packets:
$7.5\times10^6\times820\times8 = 49$\,Gbit/s (on the wire). For 100\,Gbit/s at 800 bytes: $100/49 = 2.03$, i.e.\ 3 cores
(more with headroom).
\end{solutionbox}

\begin{solutionbox}{\ref{exr:lag}}
Capacity $n\mu = 40\,000$\,s$^{-1}$. VM: headroom $200\times120 = 24\,000$, so $\rho \le 0.4$; container: $200\times10 = 2000$, so
$\rho \le 0.95$. Other limits: detection delay of the monitoring loop, image download, warm-up (caches, state
synchronisation), load-balancer reconfiguration and the availability of free hardware at the site.
\end{solutionbox}

\begin{solutionbox}{\ref{exr:pack}}
Usable capacity is 14 cores per server; total demand 32, so at least 3 servers. Placement: server 1: A (6) + B (5) + 3 = 14;
server 2: A (6) + B (5) + 3 = 14; server 3: 4 + 4 + 2 = 10. Anti-affinity holds (each server has one replica of A and B at most)
and 2 cores remain free on each: three servers.
\end{solutionbox}

% ---------------------------------------------------------------------
\section{Further reading}
The NFV architectural framework is ETSI GS NFV 002 \cite{etsinfv002} and management and orchestration ETSI GS NFV-MAN 001
\cite{etsinfvman001}. Kreutz et al.\ \cite{kreutz2015sdn} discuss the relationship between SDN and NFV; Rommer et al.\
\cite{rommer2020core} describe the cloud-native 5G core; Kleinrock \cite{kleinrock1975queueing} covers the queueing models used
for dimensioning.

\printbibliography[heading=subbibliography,title={\ACNstrRefs}]
\end{document}
